% from ega2-1-fr.tex
\oldpage[II]{28}
$\mathfrak{j}_+(Y_1)\cap\mathfrak{j}_+(Y_2)$, et à l'intersection d'une
famille quelconque $(Y_\lambda)$ de parties fermées correspond la racine
dans $S_+$ de la somme des $\mathfrak{j}_+(Y_\lambda)$.
\end{corollary}

\begin{corollary}[2.3.13]
\label{II.2.3.13-fr}
Soit $\mathfrak{J}$ un idéal gradué dans $S_+$~; pour que
$V_+(\mathfrak{J})=\varnothing$, il faut et il suffit que tout élément de
$S_+$ ait une puissance dans $\mathfrak{J}$.
\end{corollary}

Ce dernier corollaire s'exprime aussi sous l'une des formes équivalentes~:

\begin{corollary}[2.3.14]
\label{II.2.3.14-fr}
Soit $(f_\alpha)$ une famille d'éléments homogènes de $S_+$. Pour que les
$D_+(f_\alpha)$ forment un recouvrement de
$X=\operatorname{Proj}(S)$, il faut et il suffit que tout élément de $S_+$
ait une puissance dans l'idéal engendré par les $f_\alpha$.
\end{corollary}

\begin{corollary}[2.3.15]
\label{II.2.3.15-fr}
Soient $(f_\alpha)$ une famille d'éléments homogènes de $S_+$, $f$ un
élément de $S_+$. Les relations suivantes sont équivalentes~:
\begin{enumerate}
  \item[\rm a)] $D_+(f)\subset\bigcup_\alpha D_+(f_\alpha)$~;
  \item[\rm b)] $V_+(f)\supset\bigcap_\alpha V_+(f_\alpha)$~;
  \item[\rm c)] une puissance de $f$ appartient à l'idéal engendré par les
    $f_\alpha$.
\end{enumerate}
\end{corollary}

\begin{corollary}[2.3.16]
\label{II.2.3.16-fr}
Pour que $X=\operatorname{Proj}(S)$ soit vide, il faut et il suffit que tout
élément de $S_+$ soit nilpotent.
\end{corollary}

\begin{corollary}[2.3.17]
\label{II.2.3.17-fr}
Dans la correspondance biunivoque décrite dans
(\hyperref[II.2.3.12-fr]{2.3.12}), aux parties fermées irréductibles de $X$
correspondent les idéaux gradués premiers dans $S_+$.
\end{corollary}

\begin{proof}
En effet, si $Y=Y_1\cup Y_2$, où $Y_1$ et $Y_2$ sont fermés et distincts de
$Y$, on a
\[
  \mathfrak{j}_+(Y)
  =\mathfrak{j}_+(Y_1)\cap\mathfrak{j}_+(Y_2)
\]
les idéaux $\mathfrak{j}_+(Y_1)$ et $\mathfrak{j}_+(Y_2)$ étant distincts
de $\mathfrak{j}_+(Y)$, donc $\mathfrak{j}_+(Y)$ n'est pas premier.
Inversement, si $\mathfrak{J}$ est un idéal gradué non premier de $S_+$, il
existe deux éléments $f$, $g$ de $S_+$ tels que
$f\notin\mathfrak{J}$, $g\notin\mathfrak{J}$, $fg\in\mathfrak{J}$~; alors
$V_+(f)\not\supset V_+(\mathfrak{J})$,
$V_+(g)\not\supset V_+(\mathfrak{J})$ mais
$V_+(\mathfrak{J})\subset V_+(f)\cup V_+(g)$ par
(\hyperref[II.2.3.2.3-fr]{2.3.2.3})~; on en conclut que
$V_+(\mathfrak{J})$ est réunion des ensembles fermés
$V_+(f)\cap V_+(\mathfrak{J})$ et $V_+(g)\cap V_+(\mathfrak{J})$, qui sont
distincts de $V_+(\mathfrak{J})$.
\end{proof}

\subsection{La structure de schéma sur $\operatorname{Proj}(S)$.}
\label{subsection:II.2.4-fr}

\begin{env}[2.4.1]
\label{II.2.4.1-fr}
Soient $f$, $g$ deux éléments homogènes de $S_+$~; considérons les schémas
affines $Y_f=\operatorname{Spec}(S_{(f)})$,
$Y_g=\operatorname{Spec}(S_{(g)})$ et
$Y_{fg}=\operatorname{Spec}(S_{(fg)})$. En vertu de
(\hyperref[II.2.2.2-fr]{2.2.2}), le morphisme
$w_{fg,f}=({}^a\omega_{fg,f},\widetilde{\omega}_{fg,f})$ de $Y_{fg}$ dans
$Y_f$, correspondant à l'homomorphisme canonique
$\omega_{fg,f}:S_{(f)}\to S_{(fg)}$, est une \emph{immersion ouverte}
(\hyperref[I.1.3.6-fr]{I, 1.3.6}). Au moyen de l'homéomorphisme réciproque
de $\psi_f:D_+(f)\to Y_f$ (\hyperref[II.2.3.6-fr]{2.3.6}), on peut
transporter à $D_+(f)$ la structure de schéma affine de $Y_f$~; en vertu de
la commutativité du diagramme (\hyperref[II.2.3.5.1-fr]{2.3.5.1}), le
schéma affine $D_+(fg)$ s'identifie ainsi au schéma induit sur l'ensemble
ouvert $D_+(fg)$ de l'espace sous-jacent au schéma affine $D_+(f)$. Il est
clair alors (compte tenu de (\hyperref[II.2.3.4-fr]{2.3.4})) que
$X=\operatorname{Proj}(S)$ est muni d'une unique structure de
\emph{préschéma}, dont la restriction à chaque $D_+(f)$ est le schéma
affine qui vient d'être défini. En outre~:
\end{env}

\begin{proposition}[2.4.2]
\label{II.2.4.2-fr}
Le préschéma $\operatorname{Proj}(S)$ est un schéma.
\end{proposition}

\begin{proof}
Il suffit (\hyperref[I.5.5.6-fr]{I, 5.5.6}) de montrer que, quels que soient
$f$, $g$ homogènes dans $S_+$,
$D_+(f)\cap D_+(g)=D_+(fg)$ est affine et que son anneau est engendré par
les images canoniques des anneaux de $D_+(f)$ et $D_+(g)$~; le premier
point est évident par définition, et le second résulte de
(\hyperref[II.2.2.4-fr]{2.2.4}).
\end{proof}
